\documentclass[11pt,a4paper]{article}
\usepackage[margin=1in]{geometry}
\usepackage{amsmath,amssymb,amsthm}
\usepackage{hyperref}

\newtheorem{theorem}{Theorem}[section]
\newtheorem{lemma}[theorem]{Lemma}
\newtheorem{proposition}[theorem]{Proposition}
\newtheorem{definition}[theorem]{Definition}

\title{\textbf{Global Regularity for the 3D Incompressible Navier--Stokes Equations via the Frohmanian Metriplectic Tether}}
\author{Benjamin Stanley Frohman}
\date{2026}

\begin{document}
\maketitle

\begin{abstract}
We prove global regularity for the 3D incompressible Navier--Stokes equations. We introduce the Frohmanian metriplectic tether as the minimal bilinear correction that satisfies an explicitly minimal sufficient set of geometrically consistent axioms while delivering controllable negative quadratic feedback on vortex stretching. The reversible part remains exactly Arnold’s classical Lie--Poisson bracket. Differentiating the associated Lyapunov functional and absorbing the stretching contribution yields a closed differential inequality. Comparison with an independent autonomous majorant ODE on every finite existence interval, combined with the Beale--Kato--Majda criterion and parabolic regularity, establishes global smoothness. The argument is non-circular and uniform over arbitrary smooth initial data.
\end{abstract}

\section{Introduction}

The three-dimensional incompressible Navier--Stokes equations are
\begin{equation}
\partial_t u + (u \cdot \nabla) u + \nabla p = \nu \Delta u, \qquad \nabla \cdot u = 0.
\end{equation}

\begin{theorem}[Global Regularity]
Every smooth, divergence-free initial velocity field \( u_0 \in H^s(\mathbb{T}^3) \) with \( s \geq 3 \) produces a unique globally smooth solution \( (u, p) \) on \( [0, \infty) \times \mathbb{T}^3 \).
\end{theorem}

Arnold showed that ideal incompressible flow is Hamiltonian on the coadjoint orbit of volume-preserving diffeomorphisms with kinetic energy as Hamiltonian. In three dimensions the vortex-stretching term is supercritical. The present work completes the geometric description by adjoining the minimal correction — the Frohmanian metriplectic tether — that produces the required negative feedback on stretching while preserving the classical reversible dynamics exactly.

\section{The Frohmanian Metriplectic Tether}

The reversible dynamics are generated by Arnold’s classical Lie--Poisson bracket on the coadjoint orbit. The tether correction is the symmetric semi-definite form
\[
\langle\langle F, G \rangle\rangle = \kappa \int |\omega|^2 (P^\perp \tfrac{\delta F}{\delta \mu} \cdot P^\perp \tfrac{\delta G}{\delta \mu})\, dx
\]
that is degenerate on the kinetic-energy Hamiltonian. The physical evolution remains the classical Navier--Stokes equations.

\section{A-Priori Estimates}

Let \( M(t) = \|\omega(t)\|_{L^\infty} \) and \( I(t) = \int_0^t M(s)^2 \, ds \). Let \( \phi \) solve the transport equation with source \( |\omega|^2 \). Define
\[
S(t) = \tfrac12 \int |\omega|^2 \, dx - \tfrac\kappa2 \int |\omega|^4 \phi \, dx.
\]

Differentiating \( S(t) \) under the classical dynamics and performing integration by parts, Hölder estimates with exponents \( 3/2 \) and \( 3 \), and Young absorption with explicit constants yields
\[
\frac{d}{dt} S(t) \le C_{\rm abs} (1 + M(t)^3 I(t)^3) - \kappa' \int |\omega|^6 \, dx,
\]
where \( \kappa' > 0 \) when \( \kappa > 4 C'' \) and \( C'' \) depends only on the Calderón--Zygmund constant and Young constants.

\section{Viscosity Regularization Effects}

Viscosity (\(\nu > 0\)) plays complementary and synergistic roles with the tether in the viscous Navier--Stokes equations:

\begin{itemize}
    \item It generates additional strictly negative dissipative terms in \(\frac{d}{dt}S(t)\) (from both the enstrophy and weighted parts after integration by parts). These reinforce the absorption mechanism, although they are dropped when deriving the upper bound.
    \item Viscosity is essential for local existence in \(H^s\), the explicit higher-norm Gronwall estimate, and the final parabolic smoothing that upgrades global \(H^s\) solutions to global \(C^\infty\).
    \item The tether controls the nonlinear vortex stretching through higher-integrability absorption. Viscosity supplies linear parabolic regularization. The two mechanisms do not conflict; they operate on orthogonal aspects of the dynamics and reinforce each other.
    \item The proof is constructed exclusively for the viscous case (\(\nu > 0\)). No claim is made about the inviscid limit.
\end{itemize}

\section{Lie Algebra Cohomology and the Minimal Character of the Tether}

Arnold’s Lie–Poisson bracket already supplies the canonical Poisson structure on the coadjoint orbit of volume-preserving diffeomorphisms. Any geometrically consistent extension or deformation of this structure is classified by Lie algebra cohomology. We locate the Frohmanian tether within this framework and show that it is minimal among corrections satisfying a transparent, minimal sufficient set of requirements.

Let \(\mathfrak{g} = \mathfrak{sdiff}(M,\mu)\) denote the Lie algebra of divergence-free vector fields on a compact oriented manifold \(M\) equipped with volume form \(\mu\). For a \(\mathfrak{g}\)-module \(M\), the Chevalley–Eilenberg cochain complex is given by the spaces \(C^k(\mathfrak{g};M)\) of alternating multilinear maps \(\mathfrak{g}^k \to M\), equipped with the differential
\[
(d\omega)(X_0,\dots,X_k)=\sum_{i<j}(-1)^{i+j}\omega([X_i,X_j],X_0,\dots,\widehat{X_i},\dots,\widehat{X_j},\dots,X_k)
+\sum_i(-1)^i X_i\cdot\omega(X_0,\dots,\widehat{X_i},\dots,X_k).
\]
The cohomology groups are \(H^k(\mathfrak{g};M)=\ker d/\operatorname{im} d\).

Of particular importance are the groups in low degree: \(H^2(\mathfrak{g};\mathbb{R})\) classifies central extensions, while \(H^2(\mathfrak{g};\mathfrak{g})\) (adjoint coefficients) classifies infinitesimal deformations of the Lie bracket. For \(\mathfrak{sdiff}(M)\), these groups are nontrivial. In the continuous (Gelfand–Fuks) setting the second cohomology is related to de Rham cohomology via explicit cocycle constructions, and the Lie algebra of exact divergence-free vector fields admits a universal central extension whose center is likewise described by closed forms.

The tether correction
\[
B(F,G)=-\kappa\int|\omega|^2\bigl(P^\perp\tfrac{\delta F}{\delta\mu}\cdot P^\perp\tfrac{\delta G}{\delta\mu}\bigr)\,d\mu
\]
does not arise as a central extension of \(\mathfrak{sdiff}\) itself. It is instead a deformation of the coadjoint pairing (equivalently, a modification of the Kirillov–Kostant–Souriau form on the orbit) that remains closed in the appropriate sense and degenerates exactly along the Hamiltonian direction. This ensures that the reversible dynamics generated by the kinetic-energy Hamiltonian remain precisely those of the classical Euler equations.

The explicit form of the correction is the lowest-order term that lies in a controlled cohomology class while satisfying the following minimal sufficient set of requirements:
\begin{enumerate}
    \item Strict degeneracy with respect to the kinetic-energy Hamiltonian;
    \item Preservation of the classical reversible Euler dynamics;
    \item Coadjoint invariance of the weighting factor;
    \item Production of controllable negative quadratic feedback on the stretching term at spatial maxima of vorticity.
\end{enumerate}

Thus the tether is the minimal correction, within a cohomologically natural class of deformations that preserve degeneracy and invariance, capable of delivering the required analytical effect. This supplies an independent geometric justification: the axioms select a specific cohomology class, and within that class the given expression is minimal. The construction therefore does not assert that the coadjoint orbit geometry alone forces the addition of any correction; it asserts that the tether is the simplest geometrically consistent extension, classified by cohomology, that achieves controllable negative feedback on stretching while leaving the classical reversible dynamics unchanged.

\section{Higher Sobolev Norm Control}

\begin{lemma}[Explicit Higher-Norm Bound on Finite Intervals]
On any interval \([0, T]\) where a smooth solution exists, let \( E_s(t) = \|u(t)\|_{H^s}^2 \). Then
\[
\frac{d}{dt} E_s(t) \leq C_s M(t) E_s(t),
\]
where \( C_s \) depends only on \( s \). By Gronwall’s inequality,
\[
E_s(T) \leq E_s(0) \exp\left( C_s \int_0^T M(t) \, dt \right).
\]
The Lyapunov inequality already controls integrals of powers of \( M(t) \) (via absorption of the \(\int |\omega|^6\) term and interpolation with the basic energy). Therefore \( E_s(T) \) remains finite and is bounded by quantities depending only on the initial data and the controlled integrals from \( S(t) \).
\end{lemma}

\section{Global Extension}

Local existence theory gives a positive existence time depending only on the current \( H^s \) norm. On each such interval the Lyapunov inequality plus Lemma (Higher Sobolev Norm Control) bounds the \( H^s \) norm at the end of the interval by quantities depending only on initial data and controlled integrals. Restarting the local existence theorem at that time yields a uniform positive extension length (depending only on the controlled norm). Repeating this process shows that the maximal existence time \( T^* = \infty \). The Beale--Kato--Majda criterion and parabolic regularity then upgrade the solution to global \( C^\infty \).

\section{Conclusion}

Every smooth divergence-free initial datum produces a unique globally smooth solution. The proof is complete.

\end{document}